\subsection{Eliminating the four lines through the square of insertion}\label{sec:B4T01}
The parent is an ordinary residual board with $k\ge\delta_0n$ and a stationary point of $H_B$ with $0<p_e\le
C/k$, diagonal vacancies $h\ge\eta$, and weighted RC Laplacian gap at least $\gamma$. The constants are
fixed.
Take $k$ sufficiently large that $p<1/4$.
Set $q=1-p$, $\tau=pq/(1-4p)$, $T=\operatorname{diag}\tau$, and
\[
 S=\operatorname{diag}(0_{RC},h)+A^TTA,\quad
 v=(\mathbf1_R,-\mathbf1_C,0_D),\quad G=S^+.
\]
From Sections~\ref{sec:B3T05}--\ref{sec:B3T07} we use the uniform gap of $S$, along with
$G_{ij}=\delta_{ij}+O(k^{-1})$ and $\|G\|_\infty=O(1)$.
We account for the gauge using an RC cofactor. An unnormalized pseudodeterminant would give a different
factor.

\paragraph{The four removed lines and their squares.}
Let $J$ denote the four lines through $c$, let $I=J^c$, let $F$ denote the squares in their union, and set
$E'=E\setminus F$. We keep the original coordinates. Each retained line loses at most three squares. Write
\[
 d_\ell=\sum_{e\in F\cap\ell}p_e,\quad w=p|_{E'},\quad
 U=A_{F,I}^TT_FA_{F,I},\quad \Delta=\operatorname{diag}_{D\cap I}d,
\]
\[
 Q=S_{IJ}S_{JJ}^{-1}S_{JI},\quad P_c=S_{II}-Q,\quad E_c=Q+\Delta-U.
\]
Then the exact identities $S_{\rm child}(w)=P_c+E_c$ and $P_cv_I=E_cv_I=0$ hold. The matrix $Q$ has rank at
most four, whereas $\Delta-U$ need not have bounded rank.

For the feasible point $\bar p=w+z$ from Section~\ref{sec:B3T07}, we include the additional term given by
\[
 \Xi=A'^T\operatorname{diag}(\tau(\bar p)-\tau(w))A'
       -\operatorname{diag}_{D\cap I}(D'^Tz),\quad
 \mathcal E=E_c+\Xi,\quad S_{\rm child}(\bar p)=P_c+\mathcal E.
\]
Here $D'$ is the square--diagonal incidence matrix of the child.
For $y\perp v_I$, minimizing over the four removed coordinates gives
\[
 y^TP_cy=\min_a(a,y)^TS(a,y)\ge c_*\|y\|^2.
\]
Indeed, $\operatorname{dist}((a,y),\operatorname{span}v)\ge\|y\|$. Thus the kernel of $P_c$ is exactly
$\operatorname{span}v_I$.

\paragraph{Exact determinant.}
Choose the RC gauge coordinate in $I$ and define
\[
 W_c=P_c^{+1/2}\mathcal E P_c^{+1/2}\big|_{v_I^\perp}.
\]
The Schur complement followed by the cofactor identity for matrices with a common kernel gives
\[
 \boxed{\frac{\det S_{{\rm child},\rm gf}(\bar p)}{\det S_{{\rm parent},\rm gf}(p)}
 =\frac{\det_{v_I^\perp}(I+W_c)}{\det S_{JJ}}.}
\]
When using pseudodeterminants, the right-hand side must be multiplied by $(k-1)/k$, since
$\operatorname{pdet}S=2k\det S_{\rm gf}$ (Lemma~\ref{lem:puente}) and the child has $k-1$ rows.
This factor is not included a second time when working with cofactors.

\paragraph{The same elimination in $N$.}
With a retained gauge coordinate, $N=\left(\begin{smallmatrix}T^{-1}&A_{\rm gf}\\
A_{\rm gf}^T&-\operatorname{diag}h_{\rm gf}\end{smallmatrix}\right)$.
The eliminated block indexed by $(F,J)$ and the indicated block of its inverse are
\[
 K_N=\begin{pmatrix}T_F^{-1}&A_{F,J}\\A_{F,J}^T&-\operatorname{diag}h_J\end{pmatrix},
 \quad \det K_N=\Big(\prod_F\tau_e^{-1}\Big)\det S_{JJ},
\]
\[
 (K_N^{-1})_{FF}=T_F-T_FA_{F,J}S_{JJ}^{-1}A_{F,J}^TT_F.
\]
The retained Schur complement has lower block $-\operatorname{diag}h_I-U+Q$.
To obtain $N_{\rm child}(w)$, subtract $E_c$ from this lower block. For $\bar p$, also change $T_{E'}^{-1}$
and the vacancies.
This recovers $\mathfrak D=\prod_e[(1-4p_e)/q_e^4]\det S_{\rm gf}$ (Lemma~\ref{lem:puente}). We have used no
Taylor expansion taking a removed coordinate to zero or one.

\paragraph{Uniform matrix remainder.}
The incidence geometry and Section~\ref{sec:B3T07} give $\|\mathcal E\|_2=O(k^{-1})$ and $\|\mathcal
E\|_F^2=O(k^{-1})$. Indeed, $\Delta-U$ is sparse in each row, $Q$ has rank at most four and norm
$O(k^{-1})$, and $\Xi$ has absolute row sums $O(k^{-1})$ and dimension $O(k)$.
Thus $\|W_c\|_2=O(k^{-1})$, $\operatorname{tr}W_c^2=O(k^{-1})$, and
\[
 \boxed{\log\det(I+W_c)=\operatorname{tr}W_c
       -\tfrac12\operatorname{tr}W_c^2+O(k^{-2}).}
\]
The remainder is bounded by $C\|W_c\|_2\operatorname{tr}W_c^2$. The quadratic term cannot be absorbed into
an $O(k^{-2})$ remainder.
